\chapter{Le corpus : ce dont le modèle apprend}

Voici le chapitre le plus important du cours. Un entraînement parfaitement
conduit sur un mauvais corpus donne un mauvais modèle ; il n'existe aucun
réglage pour rattraper cela.

\section{Le format attendu}

Un corpus est un simple fichier texte. Chaque exemple tient sur deux lignes :
une question, une réponse.

\begin{nkcode}[]{corpus.txt — le format, en entier}
Question: Quelle est la capitale du Cameroun ?
Reponse: La capitale politique du Cameroun est Yaoundé.

Question: Comment convertir des degrés Celsius en Fahrenheit ?
Reponse: On multiplie par 1,8 puis on ajoute 32.

Question: Qu'est-ce qu'un syllogisme ?
Reponse: Un raisonnement à trois propositions où la conclusion découle
des deux premières.
\end{nkcode}

Trois règles, et c'est tout :

\begin{itemize}
  \item La ligne de question commence exactement par \texttt{Question:{} }
        (avec l'espace).
  \item La ligne de réponse commence exactement par \texttt{Reponse:{} } —
        \textbf{sans accent}, c'est le format historique du dépôt.
  \item Les lignes vides sont ignorées : mettez-en pour aérer.
\end{itemize}

\begin{nkpiege}
Une question sans réponse est \emph{silencieusement ignorée}. Si votre corpus
de 10\,000~exemples n'en donne que 4\,000 au chargement, cherchez les
\texttt{Reponse:} mal orthographiés — un accent de trop suffit.
\end{nkpiege}

\section{Combien d'exemples ?}

\begin{center}
\begin{tabular}{lp{9.2cm}}
\toprule
\textbf{Taille} & \textbf{Ce que vous pouvez en attendre} \\
\midrule
100 -- 500 & Change le \textbf{style} des réponses : leur longueur, leur ton,
             leur mise en forme. Ne change pas les connaissances. \\
1\,000 -- 10\,000 & Le modèle assimile un vocabulaire et des tournures propres à
             votre domaine. \\
50\,000 et plus & Le modèle intègre réellement le contenu de votre domaine. \\
\bottomrule
\end{tabular}
\end{center}

Le corpus utilisé pour les mesures de ce cours contient
\textbf{100\,017 paires} et pèse 25~Mo.

\begin{nkretenir}
Peu d'exemples changent la \emph{forme}, beaucoup d'exemples changent le
\emph{fond}. Si votre modèle « répond mieux » après 200 pas, il a surtout appris
à mieux se présenter — ce qui est déjà utile, mais ne le confondez pas avec du
savoir.
\end{nkretenir}

\section{La qualité prime sur la quantité}

\subsection{Ce qui fait un bon exemple}

\begin{itemize}
  \item \textbf{Une réponse exacte.} Le modèle ne vérifie rien : il imite. Une
        erreur dans le corpus devient une erreur apprise.
  \item \textbf{Une réponse complète.} « Oui » n'apprend rien. « Oui, parce
        que… » apprend à justifier.
  \item \textbf{Une formulation naturelle.} Écrivez la réponse comme vous
        voudriez la lire.
  \item \textbf{De la variété dans les questions.} Si vos mille questions
        commencent toutes par « Qu'est-ce que », le modèle ne saura répondre
        qu'à celles-là.
\end{itemize}

\subsection{Le piège du corpus fabriqué par une autre IA}

Il est tentant de demander à un modèle existant de générer des milliers de
paires question/réponse. C'est rapide, et c'est un piège.

Un modèle qui apprend des réponses d'un autre modèle apprend surtout son
\textbf{style} : sa façon de tourner les phrases, ses tics, ses formules. Il
n'apprend pas de connaissances nouvelles — l'autre modèle ne connaissait pas
votre domaine non plus.

\begin{nkpiege}
Pour \emph{spécialiser} un modèle sur votre domaine, le corpus doit venir de
\textbf{vos sources réelles} : votre documentation, vos archives, vos manuels.
Un corpus généré par une IA transmet un style, pas un savoir.
\end{nkpiege}

\section{Construire un corpus depuis votre documentation}

La méthode qui marche, en trois temps :

\begin{enumerate}
  \item \textbf{Découper.} Prenez vos documents et coupez-les en passages qui
        se suffisent à eux-mêmes — un paragraphe, une section courte.
  \item \textbf{Interroger.} Pour chaque passage, écrivez la ou les questions
        auxquelles il répond. C'est le travail le plus long, et le plus utile.
  \item \textbf{Répondre.} La réponse est le passage lui-même, reformulé pour
        qu'il se tienne seul.
\end{enumerate}

\begin{nkcode}[]{Un passage de documentation devenu deux exemples}
Question: Comment reprendre un entraînement interrompu avec NKQwen2Train ?
Reponse: Relancez exactement la même commande. Le programme cherche les
checkpoints, choisit le plus avancé qui soit intact, et reprend au pas
et à l'exemple où il s'était arrêté.

Question: Que se passe-t-il si un checkpoint est corrompu ?
Reponse: Il est détecté par son empreinte, ignoré, et la reprise se fait
depuis le second checkpoint. C'est la raison pour laquelle deux fichiers
sont conservés en alternance.
\end{nkcode}

\begin{nkretenir}
Écrire les questions est le vrai travail. C'est aussi ce qui rend le corpus
irremplaçable : il contient votre manière d'interroger votre domaine, et aucune
IA ne peut la deviner à votre place.
\end{nkretenir}

\section{Séparer l'apprentissage du contrôle}

Voici une erreur qui trompe beaucoup de débutants, et qui a une parade simple.

Si vous mesurez les progrès du modèle sur les exemples qu'il vient
d'apprendre, vous mesurez sa \textbf{mémoire}, pas sa compréhension. Un élève à
qui l'on repose exactement les questions de l'exercice qu'il vient de corriger
aura toujours l'air brillant.

La parade : mettez de côté quelques dizaines d'exemples que le modèle ne verra
\textbf{jamais} pendant l'entraînement. Ils serviront à vérifier.

\begin{itemize}
  \item Si la perte d'apprentissage descend \emph{et} celle de contrôle
        descend : le modèle apprend vraiment. \textbf{C'est ce qu'on veut.}
  \item Si la perte d'apprentissage descend mais que celle de contrôle stagne
        ou remonte : le modèle récite. On appelle cela le
        \textbf{surapprentissage}.
\end{itemize}

\texttt{NKQwen2Train} s'en charge : il évalue périodiquement un exemple pris
loin de la position courante dans le corpus, et affiche sa perte à part.

\begin{nkretenir}
Deux pertes, pas une. Celle qui compte est celle mesurée sur ce que le modèle
n'a jamais vu.
\end{nkretenir}

\section{Vérifier son corpus avant de lancer}

Un entraînement de dix jours sur un corpus mal formé est dix jours perdus.
Trois minutes de vérification les économisent.

\begin{nkcode}[]{Compter les questions et les réponses}
# Linux, macOS
grep -c "^Question: " corpus.txt
grep -c "^Reponse: "  corpus.txt

# Windows -- PowerShell
(Select-String -Path corpus.txt -Pattern '^Question: ').Count
(Select-String -Path corpus.txt -Pattern '^Reponse: ').Count
\end{nkcode}

Les deux nombres doivent être \textbf{égaux}, ou presque. Un écart important
signale des paires incomplètes.

Puis lancez l'entraînement sur deux pas seulement, juste pour lire le compte :

\begin{nkcode}[]{Vérification en trente secondes}
NKQwen2Train --gguf=<modele> --corpus=corpus.txt --steps=2 --fresh
\end{nkcode}

La ligne \texttt{corpus : N paires} vous dit combien d'exemples ont réellement
été retenus. Si ce nombre vous surprend, corrigez le corpus avant d'aller plus
loin.

\begin{nkexo}
Constituez un corpus de vingt paires sur un sujet que vous connaissez bien —
votre métier, une passion, votre ville. Écrivez-les à la main : vingt paires,
c'est une demi-heure. Vérifiez ensuite avec les commandes ci-dessus que les
vingt sont bien lues. Ce corpus vous servira aux chapitres suivants.
\end{nkexo}
